\section*{Chapter \ref{ch:ll:memory-hierarchy-and-caches} --- Memory Hierarchy and Caches}

\begin{solution}{exo:ll:memory-hierarchy-and-caches:1}
The offset is the low 6 bits: $\texttt{0x48} \bmod 64 = 8$. The set is bits 6 to 11: $\lfloor\texttt{0x5e48}/64\rfloor
\bmod 64 = 377 \bmod 64 = 57$.
\end{solution}

\begin{solution}{exo:ll:memory-hierarchy-and-caches:2}
$2\,097\,152 / (64 \times 16) = 2\,048$ sets; addresses $2\,048 \times 64 = 131\,072$ bytes (128~KiB) apart share a set.
\end{solution}

\begin{solution}{exo:ll:memory-hierarchy-and-caches:3}
$600~\text{MiB}/4~\text{KiB} = 153\,600$ entries with base pages, far beyond any TLB; 300 with 2~MiB pages.
\end{solution}

\begin{solution}{exo:ll:memory-hierarchy-and-caches:4}
\omterm{def:ll:memory-hierarchy-and-caches:false}{False sharing}: every write by one thread invalidates the other's copy of the line, and both pay a transfer between cores
on most writes (about five times the cost of an uncontended atomic increment in
\cref{fig:ll:memory-hierarchy-and-caches:false}). Put each field on its own line: 128 bytes instead of 16.
\end{solution}

\begin{solution}{exo:ll:memory-hierarchy-and-caches:5}
About $4 \times 3.5 = \qty{14}{\nano\second}$ from the second level and $4 \times 150 = \qty{600}{\nano\second}$ from
memory: the dependent loads cannot overlap.
\end{solution}

\begin{solution}{exo:ll:memory-hierarchy-and-caches:6}
In address order the next line is predictable, so the prefetcher fetches it before it is needed and the loads overlap;
in random order each address is known only when the previous load returns. Pointer-linked structures (lists, trees,
node-based maps) put their nodes wherever the allocator did, so traversals are random chases: on a hot path, prefer
contiguous arrays.
\end{solution}

\begin{solution}{exo:ll:memory-hierarchy-and-caches:7}
The bounds check is a comparison and a branch that is always predicted not taken, independent of the load's result, so
the core executes it while the dependent load it guards is outstanding. On a 64~MiB ring each step waits for memory
(about \qty{150}{\nano\second} here); a check of a cycle or two cannot show. It would show on a ring that fits the first
level, where the load itself takes five cycles.
\end{solution}

\begin{solution}{exo:ll:memory-hierarchy-and-caches:8}
One instrument and the same few messages keep the whole \omterm{def:ll:memory-hierarchy-and-caches:miss}{working set} in the first level and the predictor trained: the
\qty{20}{\nano\second} is the first-level figure. In production the updates hit many instruments and many price levels,
most of them out of the first level and some out of the second. Replay a recorded day across all instruments and
report percentiles.
\end{solution}

\begin{solution}{pb:ll:memory-hierarchy-and-caches:1}
\begin{enumerate}
\item About \qty{1.1}{\nano\second}, \qty{3.5}{\nano\second} and 100 to \qty{175}{\nano\second}.
\item At about 48~KiB, at one and a half to two megabytes, and between eight and twelve megabytes: a 48~KiB first level,
a second level of about 2~MiB, and a third level that the chase uses up to about eight megabytes.
\item The 24~MiB are shared by all the cores and hold their data too, and a random chase through the virtual machine's
pages maps lines to sets unevenly, so some sets overflow long before the whole cache is full.
\item The third level's latency, 75 cycles or more (about \qty{20}{\nano\second}); measured, 15 to
\qty{23}{\nano\second}.
\item $2 \times 256 \times 32 + 4\,096 = 20\,480$ bytes.
\item 320 lines.
\item $0.75 \times 2\,097\,152 / 20\,480 = 76.8$: 76 instruments.
\item Two, at most, with nothing else in the first level.
\item $1 - 76/300 = 74.7\%$.
\item $3 \times (0.253 \times 3.5 + 0.747 \times 150) \approx \qty{339}{\nano\second}$.
\item The 300 books span about 6~MB, 1\,500 base pages, more than the first-level TLB covers; with 2~MiB pages three
entries suffice. The gain is larger on a virtualised host, where each walk is two-dimensional.
\item Store levels closer to the touch only (a narrower ladder, chapter~19) and use smaller fields (32-bit prices and
quantities in ticks and lots); keep cold state in another structure.
\item \textbf{Named result.} On this laptop: about \qty{1.1}{\nano\second} from the first level,
\qty{3.5}{\nano\second} from the second, 100 to \qty{175}{\nano\second} from memory; 76 instruments' books fit in
three quarters of a 2~MiB second level.
\item 32 divides 64, so two levels share a line and none straddles two; with 48 bytes half the levels straddle a line
boundary and cost two misses.
\item Four arrays of one field (structure of arrays) when a scan reads one field over many levels (a depth sum); one
structure when an update reads all fields of one level.
\item The latency of each access alone, without prefetching or overlap, which a sequential benchmark hides.
\item So that the time per step is set by the buffer's size and not by a shorter cycle that fits a smaller level.
\item Several times an uncontended update, some tens of nanoseconds, each time the other thread has written it.
\item Measure the update's latency percentiles under a replay against the second-level latency, and, where hardware
counters are available, count second-level misses.
\item Where its bytes are: how many dependent loads an update makes and which level each is served from.
\end{enumerate}
\end{solution}

\begin{solution}{iq:ll:memory-hierarchy-and-caches:1}
A vector is contiguous: its elements share \omterm{def:ll:memory-hierarchy-and-caches:line}{cache lines} and the prefetcher streams them. A list's nodes are scattered, one
dependent load each, each a potential \omterm{def:ll:memory-hierarchy-and-caches:miss}{cache miss}.
\iqlookfor{contiguity, prefetching, and dependent loads.}
\end{solution}

\begin{solution}{iq:ll:memory-hierarchy-and-caches:2}
Two threads writing different variables in the same \omterm{def:ll:memory-hierarchy-and-caches:line}{cache line}, so that the line bounces between their cores. Detect it
with a profiler's cache-contention report (\texttt{perf c2c}) or by timing with padding; fix it by aligning
independently written variables to separate lines.
\iqlookfor{the coherence mechanism and padding or regrouping by writer.}
\end{solution}

\begin{solution}{iq:ll:memory-hierarchy-and-caches:3}
The cache of virtual-to-physical translations; a miss costs a page walk of several dependent loads, more under
virtualisation. \omterm{def:ll:memory-hierarchy-and-caches:tlb}{Huge pages} let one entry cover 2~MiB or 1~GiB, so the hot data's translations stay cached.
\iqlookfor{reach = entries times page size, and the page walk.}
\end{solution}

\begin{solution}{iq:ll:memory-hierarchy-and-caches:4}
A pointer chase: a random single cycle through a buffer, one slot per line, timing dependent loads for sizes from a few
kilobytes to beyond the last level; the time per step against size is a staircase. Pin the thread, use \omterm{def:ll:memory-hierarchy-and-caches:tlb}{huge pages} to
separate TLB effects, repeat.
\iqlookfor{dependent random loads defeating prefetching and overlap.}
\end{solution}

\begin{solution}{iq:ll:memory-hierarchy-and-caches:5}
Arrays of 32~KiB each, allocated page-aligned, so element $i$ of both arrays maps to the same cache set (4~KiB
aliasing): with more such streams than ways, lines evict each other. Offset one array by a line or two.
\iqlookfor{set conflicts from power-of-two strides.}
\end{solution}

\begin{solution}{iq:ll:memory-hierarchy-and-caches:6}
An open-addressing hash table with linear probing, keys and values in one flat array of small slots (order id and a
32-bit index into an order pool), sized to a power of two with a low load factor, backed by \omterm{def:ll:memory-hierarchy-and-caches:tlb}{huge pages} and
pre-faulted; one expected \omterm{def:ll:memory-hierarchy-and-caches:miss}{cache miss} per lookup, no allocation, no pointer chasing. Chapter~19 builds it.
\iqlookfor{flat storage, one miss per lookup, no allocation on insert.}
\end{solution}
